\documentclass[11pt]{beamer}

% ============================================================
% Theme
% ============================================================
\usetheme{Madrid}
\usecolortheme{seahorse}
\setbeamertemplate{navigation symbols}{}

% ============================================================
% Encoding
% ============================================================
\usepackage[utf8]{inputenc}
\usepackage[T1]{fontenc}
\usepackage{lmodern}

% ============================================================
% Packages
% ============================================================
\usepackage{amsmath}
\usepackage{listings}
\usepackage{xcolor}

\lstset{
	basicstyle=\ttfamily\small,
	frame=single,
	breaklines=true,
	columns=fullflexible,
	keywordstyle=\color{blue},
	commentstyle=\color{gray},
	stringstyle=\color{red},
	showstringspaces=false
}

% ============================================================
% Title
% ============================================================
\title{Concurrency \& Distributed Systems}
\subtitle{Part 2 --- Scaling, Limits, and Engineering Judgment}
\author{Dr. Mohamad Aoude}
\institute{Lebanese University \\ Faculty of Engineering}
\date{\today}

\begin{document}
	
	% ============================================================
	\begin{frame}
		\titlepage
	\end{frame}
	
	% ============================================================
	\section{Step 1: MultiThreadedServer}
	% ============================================================
	
	\begin{frame}[fragile]{From Sequential to Concurrent}
		\textbf{Single-threaded problem:}
		\begin{itemize}
			\item One request blocks all others
			\item Latency grows linearly under burst load
		\end{itemize}
		

		\textbf{Solution: Thread Pool}
		
		\begin{lstlisting}[language=Java]
			ExecutorService pool =
			Executors.newFixedThreadPool(10);
			
			pool.execute(() -> handleRequest(client));
		\end{lstlisting}
		

		
		\begin{block}{What you should observe}
			Latency collapses under load. Throughput increases. CPU usage increases.
		\end{block}
	\end{frame}
	
	% ------------------------------------------------------------
	\begin{frame}[fragile]{\textbf{STUDENT TASK 1} --- Run the Multi-Threaded Server}
		\textbf{Goal:} Compare single-thread vs thread-pool server.
		
		\textbf{Terminal A (server):}
		\begin{lstlisting}
			javac MultiThreadedServer.java
			java MultiThreadedServer 10
		\end{lstlisting}
		

		\textbf{Terminal B (client):}
		\begin{lstlisting}
			javac LoadTestClient.java
			java LoadTestClient --clients=50 --requests=20 --warmup=20 --csv=pool10.csv
		\end{lstlisting}

		
		\textbf{Record:} p50, p95, throughput (req/s).
	\end{frame}
	
	% ------------------------------------------------------------
	\begin{frame}[fragile]{\textbf{STUDENT TASK 2} --- Sweep Pool Size}
		\textbf{Run the same load with different pool sizes.}
		
		\begin{lstlisting}
			java MultiThreadedServer 5
			java LoadTestClient --clients=50 --requests=20 --warmup=20 --csv=pool5.csv
			
			java MultiThreadedServer 10
			java LoadTestClient --clients=50 --requests=20 --warmup=20 --csv=pool10.csv
			
			java MultiThreadedServer 50
			java LoadTestClient --clients=50 --requests=20 --warmup=20 --csv=pool50.csv
		\end{lstlisting}
		

		\textbf{Question:} Does bigger pool always improve p95?
	\end{frame}
	
	% ============================================================
	\section{Step 2: Measuring Reality}
	% ============================================================
	
	\begin{frame}{LoadTestClient --- Turning Lecture into Data}
		We measure:
		
		\begin{itemize}
			\item \textbf{p50 latency} (median)
			\item \textbf{p95 latency} (tail latency)
			\item \textbf{Throughput} (requests/sec)
		\end{itemize}
		

		
		\textbf{Why p95 matters:}
		\begin{itemize}
			\item Users remember the slowest requests
			\item Tail latency dominates perceived performance
		\end{itemize}

		
		\begin{block}{Key Shift}
			We no longer guess. We measure.
		\end{block}
	\end{frame}
	
	% ------------------------------------------------------------
	\begin{frame}{\textbf{Fill This Table (Required)}}
		\small
		\begin{tabular}{lccc}
			\textbf{Pool Size} & \textbf{p50 (ms)} & \textbf{p95 (ms)} & \textbf{Throughput (req/s)} \\
			\hline
			5   & \hspace{12mm} & \hspace{12mm} & \hspace{16mm} \\
			10  & \hspace{12mm} & \hspace{12mm} & \hspace{16mm} \\
			50  & \hspace{12mm} & \hspace{12mm} & \hspace{16mm} \\
		\end{tabular}
		
		
		\begin{alertblock}{Submission rule}
			Your table must be supported by your CSV files.
		\end{alertblock}
	\end{frame}
	
	% ============================================================
	\section{Step 3: The CPU-Bound Trap}
	% ============================================================
	
	\begin{frame}[fragile]{Replace I/O with CPU}
		Instead of (I/O wait):
		\begin{lstlisting}
			Thread.sleep(100);
		\end{lstlisting}
		
		We run (CPU work):
		\begin{lstlisting}
			int result = fibonacci(35);
		\end{lstlisting}

		
		\textbf{Now repeat the same load test.}
	\end{frame}
	
	% ------------------------------------------------------------
	\begin{frame}[fragile]{\textbf{STUDENT TASK 3} --- CPU-Bound Experiment}
		\textbf{Terminal A (server):}
		\begin{lstlisting}
			javac CPUBoundServer.java
			java CPUBoundServer 10 35
		\end{lstlisting}

		
		\textbf{Terminal B (client):}
		\begin{lstlisting}
			java LoadTestClient --clients=50 --requests=10 --warmup=20 --csv=cpu_pool10.csv
		\end{lstlisting}
		
		\pause
		
		Now increase threads:
		\begin{lstlisting}
			java CPUBoundServer 50 35
			java LoadTestClient --clients=50 --requests=10 --warmup=20 --csv=cpu_pool50.csv
		\end{lstlisting}

		
		\textbf{Question:} Did p95 improve or degrade? Why?
	\end{frame}
	
	% ------------------------------------------------------------
	\begin{frame}{What Happens Now?}
		\textbf{Expected observations:}
		\begin{itemize}
			\item CPU reaches 100\%
			\item Throughput plateaus (limited by cores)
			\item Large pools can \textbf{increase} p95 (context switching)
		\end{itemize}

		
		\begin{alertblock}{Core idea}
			Concurrency is not magic speed.
			CPU-bound workloads are limited by hardware.
		\end{alertblock}
	\end{frame}
	
	% ============================================================
	\section{Theoretical Connection}
	% ============================================================
	
	\begin{frame}{Amdahl's Law (Simplified)}
		\[
		S(N) = \frac{1}{(1 - P) + \frac{P}{N}}
		\]
		
		CPU-bound workload:
		\[
		P \approx 1
		\]
		

		Speedup limited by:
		\[
		S_{max} \approx \text{number of cores}
		\]
		

		\begin{block}{Conclusion}
			Threads do not create CPU power.
		\end{block}
	\end{frame}
	
	% ============================================================
	\section{Step 4: ThreadPoolTuner}
	% ============================================================
	
	\begin{frame}{Engineering Formula}
		\[
		\text{OptimalThreads}
		=
		\text{Cores}
		\times
		\left(1 + \frac{Wait}{Compute}\right)
		\]
		
		\textbf{Interpretation:}
		\begin{itemize}
			\item Heavy I/O $\rightarrow$ larger pools are useful
			\item CPU-bound $\rightarrow$ pool near core count
		\end{itemize}
	\end{frame}
	
	% ------------------------------------------------------------
	\begin{frame}[fragile]{\textbf{STUDENT TASK 4} --- Use ThreadPoolTuner}
		\textbf{Run:}
		\begin{lstlisting}
			javac ThreadPoolTuner.java
			java ThreadPoolTuner
		\end{lstlisting}
		

		
		\textbf{Write in your report:}
		\begin{itemize}
			\item Your machine core count
			\item Recommended pool size for (90/10), (50/50), (10/90)
			\item Compare recommendation with your best measured pool size
		\end{itemize}
	\end{frame}
	
	% ============================================================
	\section{Submission}
	% ============================================================
	
	\begin{frame}{What You Must Submit}
		\textbf{Deliverables (Lab 1.1):}
		\begin{itemize}
			\item A filled results table (pool size vs p50/p95/throughput)
			\item CSV files generated by LoadTestClient
			\item CPU-bound comparison: pool=10 vs pool=50 (explain p95 change)
			\item One paragraph: connect your results to Amdahl + tuning formula
		\end{itemize}
		
		\begin{alertblock}{Minimum acceptable report}
			If you have no CSV evidence, your numbers do not count.
		\end{alertblock}
	\end{frame}
	
	% ------------------------------------------------------------
	\begin{frame}{Final Question}
		\Large
		If threads become expensive and cores limit CPU work,
		

		what architecture avoids large thread pools entirely?
		

		\vspace{4mm}
		\begin{alertblock}{Next Topic: Event-Driven / Async I/O}
		\end{alertblock}
	\end{frame}
	
\end{document}
